\documentclass[11pt,a4paper]{article}

\usepackage[T1]{fontenc}
\usepackage{lmodern}
\usepackage{amsmath,amssymb}
\usepackage{graphicx}
\usepackage{booktabs}
\usepackage{hyperref}
\usepackage[margin=2.5cm]{geometry}

\title{Unfolding (Inverse Problems) with Uncertainty Propagation}
\author{\textit{Project write-up}}
\date{\today}

\begin{document}
\maketitle

\section{Motivation}
This short report accompanies a small learning project on inverse problems.
The goal is to estimate an underlying ``true'' binned distribution $x$ from a noisy, smeared observation $y$.
A simplified model is
\begin{equation}
  y = R x + \epsilon,
\end{equation}
where $R$ is a response matrix encoding smearing and acceptance, and $\epsilon$ represents statistical noise (often close to Poisson counting noise).

The inverse problem is typically ill-posed: small fluctuations in $y$ can cause large fluctuations in a naive inverse.
This project explores how regularization stabilizes the solution and how uncertainties can be propagated and validated.

\section{Regularized estimators}

\subsection{Tikhonov unfolding}
We estimate $x$ by minimizing a penalized least-squares objective,
\begin{equation}
  \hat{x}(\lambda) = \arg\min_x \left[ (y - Rx)^T W (y - Rx) + \lambda^2 \lVert L x \rVert^2 \right],
\end{equation}
with weight matrix $W \approx \mathrm{Cov}(y)^{-1}$ and a regularization operator $L$ (often a first- or second-difference matrix).
This gives the linear system
\begin{equation}
  (R^T W R + \lambda^2 L^T L)\,\hat{x} = R^T W y.
\end{equation}

\subsection{Truncated SVD (TSVD)}
Using $R = U S V^T$, the pseudo-inverse is $R^+ = V S^{-1} U^T$. Truncated SVD keeps only the first $k$ singular modes,
reducing variance at the cost of bias.

\section{Uncertainty propagation}
For linear estimators $\hat{x} = G y$, covariance propagation is
\begin{equation}
  \mathrm{Cov}(\hat{x}) = G\,\mathrm{Cov}(y)\,G^T.
\end{equation}
In practice we often approximate $\mathrm{Cov}(y)$ by a diagonal Poisson form $\mathrm{diag}(\max(y_i, 1))$.

\section{Validation with toy Monte Carlo}
We validate coverage and stability by generating repeated pseudo-experiments $y^{(t)} \sim \mathrm{Pois}(R x_{\mathrm{true}})$,
computing $\hat{x}^{(t)}$ and its propagated covariance, and studying pulls
\begin{equation}
  p_i = \frac{\hat{x}_i - x_{\mathrm{true},i}}{\sigma_i}, \quad \sigma_i^2 = \mathrm{Cov}(\hat{x})_{ii}.
\end{equation}
If the linearized covariance is reasonable, pull widths should be close to 1.

\section{Results}
Figures are produced by running the notebook \texttt{notebooks/demo\_unfolding.ipynb}, which saves plots into \texttt{docs/figures/}.

\begin{figure}[h]
  \centering
  \includegraphics[width=0.70\linewidth]{figures/lcurve.png}
  \caption{L-curve scan for Tikhonov unfolding (choose $\lambda$ near the corner).}
\end{figure}

\begin{figure}[h]
  \centering
  \includegraphics[width=0.95\linewidth]{figures/unfolded_vs_truth.png}
  \caption{Example unfolded solution compared to truth and the observed reco spectrum.}
\end{figure}

\begin{figure}[h]
  \centering
  \includegraphics[width=0.95\linewidth]{figures/pulls.png}
  \caption{Toy Monte Carlo pull mean and width for the chosen $\lambda$.}
\end{figure}

\section{Discussion}
Regularization controls the bias--variance tradeoff.
For this project, the most informative checks are:
(i) how unstable a naive inverse can be,
(ii) how stability changes as the regularization strength varies,
(iii) whether propagated uncertainties are consistent with toy Monte Carlo (pull widths near 1),
(iv) how sensitive conclusions are to the regularization choice.

In future iterations, I would like to expand the study to include alternative parameter-choice rules (e.g. cross-validation variants),
more realistic forward models (non-square responses, explicit efficiency), and systematic uncertainty propagation.

\end{document}
